%% kapa: exam=1 round=30 year=2019 period=1 form=A subject=civil_law subject_ko=민법 seq=9 number=9
%% source: Q-Net 감정평가사 기출문제 (cst00309) artlSeq=5205205 file=extracted/1008/5205205/2019년 제30회 감정평가사 1차시험 문제지/2019년 제30회 감정평가사 1차시험 문제지/2019년도 제30회 감정평가사 1차 1교시 A형.hwp
%% conversion: kordoc 4.24.51 -> tools/gen1.py (auto); review=proofread-sonnet; flags=none
\begin{question}{9}
물건에 관한 설명으로 옳지 않은 것은? (다툼이 있으면 판례에 따름)
\begin{choices}
  \item 관리할 수 있는 자연력은 동산이다.
  \item 주물과 종물의 법률적 운명을 달리하는 약정은 유효하다.
  \item 권원 없이 타인의 토지에서 경작한 농작물이 성숙하여 독립한 물건으로 인정되면, 그 소유권은 명인방법을 갖추지 않아도 경작자에게 있다.
  \item 특별한 사정이 없는 한, 주유기는 주유소 건물의 종물이다.
  \item 여러 개의 물건으로 이루어진 집합물은 원칙적으로 하나의 물건으로 인정된다.
\end{choices}
\end{question}
\answer{5}
